\documentclass[11pt,a4paper]{article}
\usepackage[T1]{fontenc}
\usepackage[utf8]{inputenc}
\usepackage{amsmath,amssymb,xcolor}
\usepackage[margin=25mm]{geometry}
\usepackage[colorlinks=true,urlcolor=blue]{hyperref}
\definecolor{nexorablue}{HTML}{2459B5}
\setlength{\parindent}{0pt}
\setlength{\parskip}{7pt}
\setlength{\emergencystretch}{2em}
\renewcommand{\familydefault}{\sfdefault}
\begin{document}

{\small\bfseries\color{nexorablue} NEXORA CPGE}\par
{\small MPSI / NEX-ENS-04 / Corrigé}\par
{\LARGE\bfseries Une fraction rationnelle et sa réciproque}\par
\vspace{0.25em}\hrule\vspace{0.6em}

\section*{Valeur exclue}
Si $y=2$, l’équation donne $-8=3$, impossible. Pour $y\ne2$, elle équivaut à $(y-2)x=4y+3$.

\section*{Vérification et bijection}
La solution unique vaut $(4y+3)/(y-2)$. Son écart à $4$ est $11/(y-2)$, non nul. Tout $y\ne2$ a donc un unique antécédent admissible. L’image est $\mathbb R\setminus\{2\}$ et la fraction calculée donne l’inverse.


\end{document}
